% SPDX-License-Identifier: CC-BY-SA-4.0
% Author: Matthieu Perrin
% Part: <Nom de la partie>
% Section: <Nom de la section>
% Sub-section: <Nom de la sous-section>  % (facultatif, laisser vide si non utilisé)
% Frame: <Titre de la slide>

\begingroup

\SetKwFunction{Decide}{decide}

\begin{frame}{Preuve d'indécidabilité par réduction}
  
  \begin{block}{Théorème -- Preuve d'indécidabilité par réduction}
    Soient $A$ et $B$ deux problèmes tels qu'il existe une \structure{réduction $f$ de $A$ vers $B$}. 

    Si \alert{$A$ est indécidable}, alors \alert{$B$ est indécidable}.
  \end{block}

  \begin{block}{Démonstration}
    \begin{itemize}
    \item Supposons (par contradiction) que $\structure{\Decide_B}$ décide $B$.
    \item Alors $A$ est décidé par $\structure{\Decide_A}$ :
      \begin{algorithm}[H]
        \Fun{$\Decide_A(x \in \mathcal{I}(A) ) \in \mathbb{B}$}{
          \Return $\Decide_B(f(x))$;
        }
      \end{algorithm}
    \item Absurde, car $A$ est indécidable. Donc $B$ est indécidable. 
    \end{itemize}
  \end{block}

  \pause
  
  \begin{alertblock}{Comment démontrer qu'un problème $B$ est indécidable ?}
    \begin{enumerate}
    \item Identifier un problème \alert{indécidable $A$} 
    \item Proposer une \alert{fonction algorithmique} $f : \mathcal{I}(A) \rightarrow \mathcal{I}(B)$ 
    \item Justifier que $\forall x \in \mathcal{I}(A),\quad \alert{x\in \textsc{pos}(A) \Leftrightarrow f(x) \in  \textsc{pos}(B)}$
    \end{enumerate}
  \end{alertblock}
  
\end{frame}

\endgroup
